% ===== §3 Antithesis: self-continuation is morally indifferent =====
\section{Antithesis: Self-Continuation Is Morally Indifferent}
\label{sec:antithesis}

\noindent
The thesis stands; and against it the antithesis raises a single, stubborn law. \emb{Self-continuation, in itself, is morally indifferent.} That a cycle perpetuates itself tells us nothing of its goodness, for the most rapacious cycles perpetuate themselves no less than the most generous---and some of them perpetuate themselves precisely \emph{by} their rapacity. The autotelic structure that the thesis identified as a metaphysical ground is, taken alone, no mark of the good at all; it is, if anything, the formal signature of the very thing the paper means to criticize. This is the hard turn the argument must take before it can build: the ground it has just laid does not yet bear any ethical weight, and to pretend otherwise would be to mistake a structure for a value.

Consider the cases, set deliberately side by side, for their likeness is the whole of the point.

\emc{Capital} is the purest case, and the one in which the autotelic structure is most exactly visible. Marx's formula for capital is not C--M--C---commodity exchanged, by way of money, for another commodity, to satisfy a need that ends the movement---but M--C--M$'$: money advanced to become more money, a movement whose end is not any commodity, not the satisfaction of any need, but its own continuation at an enlarged scale \citep{marx1976}. Capital, in Marx's striking phrase, is value in process, value that has become the ``automatic subject'' of its own movement, and its aim is nothing but its own valorization---accumulation for the sake of accumulation, production for the sake of more production. This is self-continuation in its most exact and most relentless form: production pointing at production, the cycle whose only end is the cycle. And it is precisely the thing the paper means to criticize. Here is the first and sharpest blow to any easy reading of the thesis: the autotelic form, far from being the mark of the good, is the formal signature of what the paper will come to call the vicious circle. M--C--M$'$ satisfies the thesis to the letter, and it is the model of extraction.

\emc{The rose} is the case carried over from the prelude, and it shows that the indifference of self-continuation is not peculiar to the economic. The rose's blooming is self-continuation of a kind---the plant's drive to flower, to set seed, to flower again, a generativity oriented toward its own perpetuation through the generations. And yet, at the scale of the single flower, it is exhaustive: the resplendence is bought by a total expenditure that holds nothing in reserve, and the bloom that spends everything is, for that reason, the bloom that cannot last. The rose perpetuates the species by spending the flower; and whether this spending is good or ill---whether it is the generous fall that feeds the root or the anxious blaze that exhausts it---the bare fact of perpetuation cannot say. Self-continuation here wears neither a virtuous nor a vicious face; it is, in itself, indifferent between them.

\emc{The parasitic bond} is the case nearest the paper's true subject, and the one that should most disturb. A relation in which one party's self-continuation is purchased by the steady consumption of the other---the bond that endures, sometimes for decades, precisely by extracting from one partner the resource that sustains the other---perpetuates itself as surely as any love, and often more stably, for the extracted party's depletion can be made to look like devotion and the arrangement thereby stabilized. It is a cycle; it does not stop; by the thesis's criterion it is a fully successful generativity. And it is vicious. That a bond lasts, that it reproduces itself across years, that value circulates within it---none of this distinguishes it from love, and the failure to see this is the failure the antithesis exists to correct.

The three are formally alike: each is a self-continuing cycle, each satisfies the thesis, and the likeness is exactly what exposes the insufficiency of the thesis---and, with it, the insufficiency of Eglash's spatial criterion taken alone.

\begin{claim}[The insufficiency of the spatial criterion]
The recursive return of value---Eglash's spatial criterion---is necessary to the good circle but does not suffice for it. For it cannot, by itself, exclude two failures. It cannot exclude the \emph{fair stasis}: a cycle in which value duly returns to its creators, and which is in this sense just, yet which merely repeats, generating no surplus, climbing nowhere---the bond grown fair and lifeless, conserved but not sublimating, a circle and not a spiral. And it cannot exclude the \emph{counterfeit closure}: a cycle that appears to return value to its creators while in fact sustaining itself by extraction from an outside it has hidden from view. Self-continuation, even self-continuation \emph{with} return, is not yet the good.
\end{claim}

The two failures deserve to be drawn out, for the rest of the paper is, in large part, the apparatus for detecting them.

\emc{The fair stasis} is the gentler failure, and the one the spatial criterion is least equipped to see, since by the spatial criterion it is not a failure at all. Imagine a bond in which value genuinely returns to both parties, neither extracts from the other, the distribution is scrupulously just---and nothing grows. The cycle turns and returns to exactly where it began; each round reproduces the last without remainder; the relation is fair, and dead. This is the circle of zero holonomy that the synthesis will name: a return that returns to the same, accruing nothing, climbing nowhere. The spatial criterion pronounces it just and falls silent, for it has no temporal dimension with which to register that a just cycle may yet be a stagnant one. That a relation can be fair and lifeless at once is the first thing the spatial criterion cannot say---and the first reason a temporal criterion is needed.

\emc{The counterfeit closure} is the graver failure, and it deserves the closer look, for it is the structural heart of the antithesis and the thing the whole paper must learn to detect. Here is the puzzle it poses: if the vicious cycle is truly unsustainable, how does it endure---often for a very long time, capital for centuries, the parasitic bond for decades? A theory of sustainability that could not answer this would be mere moral consolation, predicting a collapse that stubbornly fails to come. The answer is that \emb{the vicious cycle is an open system masquerading as a closed loop.} It sustains itself not by internally regenerating its conditions but by \emph{dumping its entropy on an outside}---a colony, a nature, a future generation, a silenced party---and it defers its own collapse by an ever-widening extraction. Its ``closure'' is space bought with time: so long as there remains a fresh outside to draw upon, the breach in the loop stays hidden, and the system can present itself as self-sustaining while in truth it is being fed from a source it does not acknowledge. This is why capital must perpetually expand---to new markets, to the financialization of nature, to the colonization of the future; expansion is not capital's ambition but its necessity, the continual search for a new outside to consume now that the last is exhausted. And it is why the parasitic bond so often goes hand in hand with the isolation of its victim, the quiet severing of the partner's external supports---so that there is nowhere for the extracted to flee, and no witness to the extraction. \emb{The vicious cycle's endurance is the postponement, and the displacement onto another, of its own collapse.}

Here, then, is the question the antithesis hands to the synthesis. If the cycle is, as the thesis and Eglash alike suggest, all but universal---if generativity must lean on recursion to perpetuate itself at all---and if self-continuation cannot tell good from ill, and if even the return of value cannot exclude the fair stasis or the counterfeit closure, then \emb{wherein lies the difference between the good circle and the vicious one}? The thesis gave us the universality of the cycle; the antithesis has given us its moral indifference, and named the two failures---the stagnant and the counterfeit---that any adequate criterion must exclude. The synthesis must now give us the criterion that the cycle, by itself, withholds: a difference not in the bare fact of self-continuation, which good and vicious share, but in the manner of the return.
